\section{Even and odd channels as diagnostic residuals}
We want to remove the legitimate symmetric magnetic structure while retaining
components that the model forbids. Compare points at equal $i_d$, temperature,
and signed speed, with opposite equal-magnitude $i_q$. Define
\begin{equation}
 \evenq[f](i_d,i_q)=\frac{f(i_d,i_q)+f(i_d,-i_q)}{2},\qquad
 \oddq[f](i_d,i_q)=\frac{f(i_d,i_q)-f(i_d,-i_q)}{2}.
 \label{eq:parity-operators}
\end{equation}
Addition retains the unchanged part; subtraction retains the reversing part.
These are complementary projections: $f=\evenq[f]+\oddq[f]$,
$\evenq^2=\evenq$, $\oddq^2=\oddq$, and $\evenq\oddq=0$. Pairing therefore discards the
allowed component in each diagnostic channel, not saturation as such.
\begin{figure}[tb]
\centering
\begin{tikzpicture}
\begin{axis}[diagnostic axis,title={Reflect the same curve},ylabel={$f(y)$},
 xmin=-1,xmax=1,ymin=.5,ymax=1.6]
 \addplot[reference quantity,strong line,domain=-1:1]{1+.2*x^2+.3*x};
 \addlegendentry{$f(y)$}
 \addplot[estimated quantity,alternative pattern,standard line,domain=-1:1]{1+.2*x^2-.3*x};
 \addlegendentry{$f(-y)$}
\end{axis}
\begin{axis}[diagnostic axis,at={(.5\textwidth,0)},anchor=south west,
 title={Average and half-difference},ylabel={Component},xmin=-1,xmax=1,ymin=-.4,ymax=1.4]
 \addplot[derived quantity,strong line,domain=-1:1]{1+.2*x^2};
 \addlegendentry{$\evenq[f]$}
 \addplot[torque quantity,standard line,domain=-1:1]{.3*x};
 \addlegendentry{$\oddq[f]$}
\end{axis}
\end{tikzpicture}
\caption{Reflection separates an illustrative curve
$f(y)=1+0.2y^2+0.3y$ into its allowed even part and its odd tilt.
The same projections applied to d and q flux select opposite diagnostic
channels. The schematic curve is not a second machine model.}
\label{fig:decomposition}
\end{figure}


For a correctly aligned ideal map, $\oddq[\psid]=0$ and $\evenq[\psiq]=0$.
Define measured residuals at the observed numerical labels by
\begin{equation}
 \boxed{r_d=\oddq[\hat\psid],\qquad r_q=\evenq[\hat\psiq].}
 \label{eq:residuals}
\end{equation}
For the independent positive-q representative $a>0$, this means explicitly
\begin{equation}
 r_d=\frac{\hat\psi_d(i_d,+a)-\hat\psi_d(i_d,-a)}{2},\qquad
 r_q=\frac{\hat\psi_q(i_d,+a)+\hat\psi_q(i_d,-a)}{2}.
\end{equation}

\subsection{Symmetry projection as a correction operator}
The same even/odd decomposition defines the closest parity-admissible pairwise map under equal weighting:
\begin{equation}
 \boxed{\psi_d^{\rm sym}=\evenq[\hat\psi_d],\qquad
        \psi_q^{\rm sym}=\oddq[\hat\psi_q].}
 \label{eq:symmetry-projection}
\end{equation}
The removed component is exactly
\begin{equation}
 \hat\flux-\flux^{\rm sym}=\begin{bmatrix}r_d\\r_q\end{bmatrix}
\end{equation}
on mirrored support. Algebraically this is an idempotent projection onto the declared parity subspace. Geometrically it folds the two measured branches onto the reflection-compatible component. Physically it is a correction only if the forbidden component is known to be nonmagnetic or otherwise inadmissible for the intended model.

A smaller symmetry residual after projection is therefore guaranteed by construction and cannot by itself validate a sensor correction, winding resistance, or rotor-angle estimate. The raw field, projected field, and removed residual must remain available separately.

\paragraph{What pairing hides.}
An allowed-parity flux error can be large and yet yield zero residuals.
A nonzero residual indicates incompatibility with the specified symmetry,
not unique proof of resistance or angle error. Interpolation mismatch,
temperature mismatch, or real machine asymmetry can enter the same channels.
Use one independent representative per pair, conventionally $i_q>0$; keeping
both representatives doubles the same information and their correlated noise.


\subsection{Why the closest pair is an average}
For two observed d values $d_+$ and $d_-$, the admissible pair is $(d,d)$.
Minimizing $(d-d_+)^2+(d-d_-)^2$ gives $d=(d_++d_-)/2$.
For q values the admissible pair is $(q,-q)$ and minimizing
$(q-q_+)^2+(-q-q_-)^2$ gives $q=(q_+-q_-)/2$.
These two one-dimensional minimizations prove the minimum-distance claim.
With unequal independent precisions $w_+,w_-$, replace the averages by
$(w_+d_++w_-d_-)/(w_++w_-)$ and
$(w_+q_+-w_-q_-)/(w_++w_-)$. Correlated uncertainty requires its full metric.
Equal weighting is a declared mathematical choice, not an inferred sensor
covariance. A normalized thought experiment $(d_+,d_-)=(3,1)$ projects to
$(2,2)$ and removes $(1,-1)$; the operator cannot tell why the difference arose.
